\documentclass[12pt,a4paper]{article}
\usepackage[utf8]{inputenc}
\usepackage[T1]{fontenc}
\usepackage{amsmath,amssymb,graphicx,booktabs,hyperref,natbib}
\usepackage[margin=2.5cm]{geometry}

\title{\bf Supplementary Information:\\
Dark energy as an information capacity deficit\\
--- Full Derivations, Extended Data, and Additional Tests}

\date{June 2026}

\begin{document}
\maketitle
\tableofcontents

\section{The DGF Framework: Complete Exposition}

\subsection{Axioms}

The Discrete Graph Framework (DGF) is built on two axioms:

\vspace{4pt}
\noindent\fbox{\begin{minipage}{0.95\textwidth}
{\bf A1---Causal existence.} There exist asymmetric influence relations: ``$A$ affects $B$'' is not equivalent to ``$B$ affects $A$.'' This is the minimal operational definition of existence---to exist is to be distinguishable, and distinguishability requires directional influence.\\
{\bf A2---Bounded capacity.} Each minimal causal unit (graph node) carries at most 1 bit of distinguishable information. The maximum propagation speed of causal influence is $c$. The global state space is therefore $(\mathbb{C}^2)^{\otimes N}\simeq\mathbb{C}^{2^N}$.
\end{minipage}}

\vspace{4pt}
From these axioms, several structural properties follow:

\begin{enumerate}
\item {\bf Finite-dimensional state space.} A2 directly implies the Hilbert space dimension is $2^N$, finite but exponentially large. No continuum limit is assumed.
\item {\bf Dynamics are permutations.} Evolution within a finite state space that preserves distinguishability (A1) must be a permutation of basis elements. Continuous time evolution emerges from the limit of many small permutation steps.
\item {\bf The causal graph.} Nodes are minimal causal units. Edges are directed influence relations (A1). The graph topology encodes the full causal structure of the universe.
\item {\bf The order parameter $q$.} Define $q\equiv N_0/N$, where $N_0$ is the number of nodes in state $|0\rangle$ (``unoccupied''---no information content beyond vacuum zero-point fluctuations) and $N$ is the total number of nodes. $q=1$ corresponds to pure vacuum; $q=0$ to complete classicalization. $\Delta\equiv1-q$ is the information deficit.
\end{enumerate}

\subsection{Theorem A3: The Causal Horizon}

{\bf Statement:} The permutation of $|0\rangle$ (unoccupied) nodes is causally indistinguishable---no observable can distinguish which specific $|0\rangle$ node received a causal influence. The identity of $|1\rangle$ nodes is fixed by their causal history.

{\bf Proof sketch:} From A1+A2, a fully occupied node ($|1\rangle$) cannot ``absorb'' further causal influence without violating the 1-bit capacity bound (A2). It must transmit influence to an unoccupied neighbor. That neighbor becomes $|1\rangle$ (determined). But which specific $|0\rangle$ node received the transmission is not recorded in any causal history---the permutation of $|0\rangle$ nodes is a symmetry of the causal DAG. $\square$

This was originally Axiom A3; it was demoted to a theorem when shown to follow from A1+A2+graph connectivity (S4).

\section{Derivation of the Telegraph Equation}

\subsection{From permutations to continuous evolution}

Consider $N$ causal graph nodes. The state of each node is $|0\rangle$ (unoccupied) or $|1\rangle$ (occupied, determined). The global state is a vector in $(\mathbb{C}^2)^{\otimes N}$.

The dynamics are a sequence of local permutations (gate operations). In the large-$N$ limit, the coarse-grained dynamics of $q=N_0/N$ obey a Fokker-Planck equation. The drift term comes from the mean rate of $|0\rangle\to|1\rangle$ transitions (driven by structure formation); the diffusion term comes from fluctuations in the occupation process.

\subsection{The nonlinear telegraph equation}

The resulting evolution equation for $q(x^\mu)$ is:

\begin{equation}\label{sm:telegraph}
\boxed{\partial_t^2 q + \gamma_0(1-q)\partial_t q = c^2\nabla^2(\ln q) - \kappa\int_0^t [1-q(t')]dt' + \eta\,\dot{\rho}_*(t)}
\end{equation}

Term-by-term physics:
\begin{itemize}
\item $\partial_t^2 q$: Inertial term. The $q$-field has effective ``mass'' from the finite speed of causal propagation.
\item $\gamma_0(1-q)\partial_t q$: {\bf Nonlinear damping.} Damping strength $\propto(1-q)=\Delta$, the information deficit. Pure vacuum ($q=1$) has zero damping; a fully occupied universe ($q=0$) has maximal damping. This is the natural consequence of occupied modes experiencing ``friction'' from their mutual causal connections.
\item $c^2\nabla^2(\ln q)$: {\bf Spatial stiffness.} The Laplacian of Shannon self-information $-\ln q$. In the static limit ($\partial_t q=0$), this yields $\nabla^2(\ln q)=0$, whose spherically symmetric solution $q(r)=e^{-GM/rc^2}$ recovers Newtonian gravity in the weak-field limit.
\item $-\kappa\int_0^t [1-q(t')]dt'$: {\bf Nonlocal memory.} The accumulated history of mode occupation acts as a restoring force---the universe ``remembers'' past structure formation and resists further deviation from $q=1$.
\item $+\eta\,\dot{\rho}_*(t)$: {\bf Driving by structure formation.} New stars and galaxies occupy vacuum information modes, driving $q$ away from 1.
\end{itemize}

\subsection{Generalized damping nonlinearity}

We generalize the damping to $\gamma_0(1-q)^n\dot{q}$ with $n\in[0,\infty)$:

\begin{equation}
\partial_t^2 q + \gamma_0(1-q)^n\partial_t q = c^2\nabla^2(\ln q) - \kappa\int_0^t [1-q(t')]dt' + \eta\,\dot{\rho}_*(t)
\end{equation}

The natural benchmark $n=1$ corresponds to damping proportional to the information deficit $\Delta$ itself. Alternative $n$ values correspond to different statistical mechanics of information-mode interactions:
\begin{itemize}
\item $n=0$: Linear damping (independent of $\Delta$)---excluded by DESI DR2 at $>6\sigma$
\item $n=1$: Damping $\propto\Delta$ (natural benchmark)---selected by DESI DR2
\item $n=2$: Damping $\propto\Delta^2$---marginally excluded
\end{itemize}

\subsection{Slow-roll approximation}

In the cosmological context, the inertial term $\partial_t^2 q$ is subdominant to the damping and driving terms. The slow-roll condition is:
\begin{equation}
|\partial_t^2 q| \ll |\gamma_0\Delta^n\partial_t q|,\; |\kappa\!\int\!\Delta dt'|,\; |\eta\dot{\rho}_*|
\end{equation}

Under this condition, and on scales where $\nabla^2(\ln q)\approx0$ (homogeneous cosmology):
\begin{equation}\label{sm:slowroll}
\boxed{\gamma_0\Delta^n\frac{d\Delta}{dt} + \kappa\int_0^t\Delta\,dt' = \eta\,\dot{\rho}_*(t)}
\end{equation}

where we have used $q=1-\Delta$, $\dot{q}=-\dot{\Delta}$.

\section{Derivation of the $w(z)$ Shape Function}

\subsection{Driving-dominated regime}

For $z\gtrsim0.5$, the memory integral term is subdominant to the SFR driving term. We verify this by computing the memory-drive ratio:
\begin{equation}
\mathcal{R}(t)\equiv\frac{\kappa\int_0^t\Delta\,dt'}{\eta\,\dot{\rho}_*(t)}
\end{equation}

At $z=0$, using calibrated parameters, $\mathcal{R}_0\approx0.3$; at $z=1$, $\mathcal{R}\approx0.1$. The driving-dominated approximation is accurate to $\sim10\%$--$30\%$.

In the driving-dominated limit, equation (\ref{sm:slowroll}) becomes:
\begin{equation}
\gamma_0\Delta^n\frac{d\Delta}{dt} = \eta\,\dot{\rho}_*(t)
\end{equation}

Multiplying by $\Delta^n$ and integrating:
\begin{equation}
\frac{\gamma_0}{n+1}\frac{d}{dt}(\Delta^{n+1}) = \eta\,\dot{\rho}_*(t) \quad\Rightarrow\quad \Delta^{n+1}(t) = \frac{(n+1)\eta}{\gamma_0}\rho_*(t)
\end{equation}

Thus:
\begin{equation}\label{sm:delta}
\boxed{\Delta(t) = \left[\frac{(n+1)\eta}{\gamma_0}\right]^{1/(n+1)} \rho_*(t)^{1/(n+1)}}
\end{equation}

\subsection{Equation of state}

From energy conservation:
\begin{equation}
\dot{\rho}_{\rm DE} + 3H(\rho_{\rm DE}+P_{\rm DE}) = 0
\end{equation}

With $\rho_{\rm DE}=\rho_\Lambda q = \rho_\Lambda(1-\Delta)$:
\begin{equation}
w = \frac{P_{\rm DE}}{\rho_{\rm DE}} = -1 - \frac{\dot{\rho}_{\rm DE}}{3H\rho_{\rm DE}} = -1 + \frac{\dot{\Delta}}{3H(1-\Delta)}
\end{equation}

Substituting $\dot{\Delta}$ from (\ref{sm:slowroll}) in the driving-dominated limit:
\begin{equation}
\dot{\Delta} = \frac{\eta\,\dot{\rho}_*}{\gamma_0\Delta^n} = \frac{\eta}{\gamma_0}\cdot\frac{{\rm SFR}(t)}{[(n+1)\eta/\gamma_0]^{n/(n+1)}\rho_*(t)^{n/(n+1)}}
\end{equation}

Therefore:
\begin{equation}\label{sm:wfinal}
\boxed{w(z)+1 = \frac{1}{3}\sqrt{\frac{\eta}{2\gamma_0}}\;\cdot\;\frac{{\rm SFR}(z)}{H(z)\,\rho_*(z)^{n/(n+1)}}\;\cdot\;\frac{1}{1-\Delta(z)}}
\end{equation}

For $\Delta\ll1$, the last factor $\approx1$, yielding the simplified form used in the main text. The full form including the $(1-\Delta)^{-1}$ correction is used in the numerical computations.

\subsection{Normalized shape function}

Defining the coupling constant $C\equiv\frac{1}{3}\sqrt{\eta/(2\gamma_0)}$, the normalized shape is:
\begin{equation}
\boxed{S(z)\equiv\frac{w(z)+1}{w(0)+1} = \frac{{\rm SFR}(z)/[H(z)\rho_*(z)^{n/(n+1)}]}{{\rm SFR}(0)/[H_0\rho_{*,0}^{n/(n+1)}]}}
\end{equation}

$S(z)$ is independent of the coupling constant $C$---it depends only on known astrophysical quantities and the single physical parameter $n$.

\section{Cross-Framework Consistency: S1, S3, S4}

\subsection{S1: Quantum mechanics from permutation dynamics}

DGF derives the full mathematical structure of quantum mechanics from A1+A2+graph indistinguishability (see LP32-S1). The key result relevant to cosmology is:

{\bf S1 Theorem} (Quantum-Classical Boundary): The probability of ``information reflux''---a determined $|1\rangle$ node returning to $|0\rangle$---is bounded above by:
\begin{equation}
P_{\rm reflux} \leq \frac{q_S}{q_E}
\end{equation}
where $q_S$ is the system's free capacity and $q_E$ is the environment's free capacity. In the cosmological context, this gives the upper bound on the memory recovery rate:
\begin{equation}
\dot{\Delta}_{\rm recovery} \leq \frac{q}{1-q}\,H_0\Delta
\end{equation}

This bound is satisfied by all our numerical solutions.

\subsection{S3: Einstein equations from $q$-field action}

DGF derives the Einstein field equations from the $q$-field effective action (see LP32-S3):
\begin{equation}
S_q = \int\sqrt{-g}\,d^4x\left[\frac{1}{2\kappa}\frac{(\partial q)^2}{q^2} + V(q) + \xi(q)R\right]
\end{equation}

Varying with respect to $g^{\mu\nu}$ yields:
\begin{equation}
G_{\mu\nu} = 8\pi G_{\rm eff}(q)\left[T_{\mu\nu}^{\rm(matter)} + T_{\mu\nu}^{(q)}\right]
\end{equation}
where $G_{\rm eff}(q)=G/[1+16\pi G\xi(q)]$. In the limit $q\to1$ (pure vacuum), standard GR is recovered. The $f(q)R$ non-minimal coupling avoids the Li-Pang no-go theorem for entropic gravity.

\subsection{S4: Black hole entropy area law}

DGF provides a rigorous proof of the Bekenstein-Hawking area law using only A1+A2 (see LP32-S4):
\begin{equation}
S \leq N_{\partial\Omega} \propto A
\end{equation}

The proof uses the Ford-Fulkerson min-cut theorem on the causal graph: the number of distinguishable external configurations of a black hole is bounded by the number of graph edges crossing its boundary. This is a {\it combinatorial} result---no quantum field theory in curved spacetime is required.

\section{Error Propagation Analysis}

\subsection{Theory error budget for $w_0$}

\begin{table}[htbp]
\centering
\caption{Error budget for $w_0+1$ in DGF $n=1$.}
\begin{tabular}{lccc}
\toprule
Source & Current uncertainty & Contribution to $\sigma_{w_0+1}$ & Fraction \\
\midrule
SFR$_0$ normalization & $\pm27\%$ \cite{madau2014,driver2018} & $\pm0.054$ & 73\% \\
$\rho_{*,0}$ normalization & $\pm11\%$ \cite{driver2018} & $\pm0.006$ & 8\% \\
$n$ (damping index) & $n\in[0.7,1.3]$ & $\pm0.020$ & 14\% \\
$H_0$ & $\pm0.7\%$ \cite{planck2018} & $\pm0.001$ & 5\% \\
\midrule
{\bf Total} & & {\bf $\pm0.058$} & 100\% \\
DESI measurement & & $\pm0.047$ & --- \\
\bottomrule
\end{tabular}
\end{table}

The theory uncertainty is currently dominated by the SFR$_0$ systematics. With multi-band SFR constraints (H$\alpha$+UV+FUV), the SFR$_0$ uncertainty could be reduced to $\pm10\%$, bringing the total theory error to $\pm0.024$---below the DESI measurement error.

\subsection{Shape function uncertainty}

The normalized shape $S(z)$ is independent of the absolute SFR and $\rho_*$ normalizations (they cancel in the ratio). Its uncertainty comes from:
\begin{enumerate}
\item The SFR {\it shape} as a function of redshift (dust extinction corrections at $z>2$)
\item The $\rho_*$ integration (choice of IMF, mass return fraction 0.72)
\item The memory kernel form (constant vs.\ frequency-dependent)
\end{enumerate}

At $z<2$, the shape uncertainty is $\sim\pm15\%$; at $z>2$, it grows to $\sim\pm30\%$ due to increasing SFR measurement uncertainties.

\section{Quantum Darwinism: Laboratory Test of $q$}

\subsection{The DGF prediction}

In quantum Darwinism \cite{zurek2009}, a system $S$ interacts with an environment $E$ of $N$ fragments. The redundancy $R_\delta$ is the number of independent environment fragments that contain (nearly) complete information about $S$.

In DGF, the available redundancy is proportional to the fraction of environment modes that are unoccupied:
\begin{equation}\label{sm:rdelta}
\boxed{R_\delta(q) = R_\delta(1)\cdot q}
\end{equation}
where $R_\delta(1)$ is the maximum redundancy when all $N$ environment modes are available.

\subsection{Experimental protocol}

\begin{enumerate}
\item Prepare $N=20$ environment qubits in a superconducting circuit
\item Pre-occupy a fraction $1-q$ of them by initializing to $|1\rangle$ (simulating DGF structure formation)
\item Prepare system qubit $S$ in $|+\rangle$, entangle with all environment qubits via CNOT
\item Measure mutual information $I(S:F_k)$ for fragments $F_k$ of size $k$
\item Determine $R_\delta(q)$ as the smallest $k$ such that $I(S:F_k)\approx H(S)=1$ bit
\item Repeat for different $q$ values and verify the linear relationship (\ref{sm:rdelta})
\end{enumerate}

\subsection{Distinguishing DGF from standard decoherence}

The key distinction: in standard quantum mechanics, pre-occupying environment qubits merely reduces the effective environment size ($N\to qN$), making $R_\delta\propto qN/N_0$. In DGF, $R_\delta\propto q$ with a {\it fixed} slope $R_\delta(1)$, independent of the system-environment coupling strength.

By varying the CNOT gate time (changing the effective coupling), one can distinguish these predictions:
\begin{itemize}
\item Standard QM: $R_\delta(q)$ slope depends on coupling strength
\item DGF: $R_\delta(q)$ slope is fixed at $R_\delta(1)$
\end{itemize}

\subsection{Feasibility}

Zhu et al.\ (2025, {\it Science Advances}) have demonstrated quantum Darwinism measurements in 10-qubit systems with $\sim15\%$ redundancy precision. For the DGF test, $\sim20$ qubits and $\sim5\%$ precision are needed---feasible within 2--3 years on IBM's 1000-qubit roadmap.

\section{Extended Competitor Comparison}

\begin{table}[htbp]
\centering
\caption{\bf Complete comparison of dark energy models.}
\begin{tabular}{lccccc}
\toprule
Model & $N_{\rm par}$ & Physical mechanism & Non-cosmology tests? & Non-monotonic $w(z)$? & $\chi^2$ (DESI DR2) \\
\midrule
$\Lambda$CDM & 0 & None & No & No & 57.5 \\
CPL & 2 & None (phenomenological) & No & No & 0.0 \\
$w_0w_a$CDM & 2 & None (phenomenological) & No & No & 0.0 \\
Quintessence ($V\propto\phi^{-\alpha}$) & 2--3 & Scalar field potential & No & No & $\sim$0--2 \\
Emergent DE \cite{li2024} & 2--3 & $\Lambda$ as integration constant & Partial & Possible & $\sim$0--2 \\
QMM \cite{neukart2025} & 2 & Quantum memory modes & In principle & No & $>100$ (excluded) \\
Holographic DE & 1--2 & Horizon entropy bound & No & No & $\sim$5--10 \\
CreSS & 3 & Curvature-driven & No & Possible & $\sim$0--3 \\
{\bf DGF (this work)} & {\bf 2} & {\bf Information capacity deficit} & {\bf Yes (S1,S4,S5)} & {\bf Yes} & {\bf 0.65} \\
\bottomrule
\end{tabular}
\end{table}

\section{Additional Figures}

[Figures 2 and 3 from the main text, plus supplementary figures showing: the memory-drive ratio $\mathcal{R}(z)$, the $\Delta(z)$ evolution, and the SFR uncertainty propagation to $S(z)$.]

\begin{thebibliography}{99}
\bibitem{madau2014} Madau, P.\ \& Dickinson, M.\ Cosmic star-formation history. {\it Annu.\ Rev.\ Astron.\ Astrophys.} {\bf 52}, 415--486 (2014).
\bibitem{driver2018} Driver, S.P.\ et al.\ GAMA/G10-COSMOS/3D-HST. {\it Mon.\ Not.\ R.\ Astron.\ Soc.} {\bf 475}, 2891--2935 (2018).
\bibitem{planck2018} Planck Collaboration. Planck 2018 results. VI. {\it Astron.\ Astrophys.} {\bf 641}, A6 (2020).
\bibitem{zurek2009} Zurek, W.H.\ Quantum Darwinism. {\it Nature Phys.} {\bf 5}, 181--188 (2009).
\bibitem{li2024} Li, X.\ et al.\ Emergent dark energy. {\it Phys.\ Rev.\ D} {\bf 109}, 043510 (2024).
\bibitem{neukart2025} Neukart, F.\ et al.\ QMM model for dark energy. {\it Astronomy} {\bf 4}, 16 (2025).
\end{thebibliography}

\end{document}
